% from ega4-12-fr.tex
\oldpage[IV-3]{186}
on voit qu'on peut, pour prouver le lemme, supposer en
outre $v$ surjectif. Le lemme sera alors conséquence du suivant~:

\begin{lemma}[12.3.3.3]
\label{IV.12.3.3.3-fr}
Soient $B$ un anneau, $M,N,P$ trois $B$-modules,
$u:M\to N$, $v:N\to P$ deux homomorphismes tels que $u$ soit injectif, $v$
surjectif et $v\circ u=0$. Soit $t$ un élément de $B$ tel que $t$ soit
$P$-régulier. Alors on a
\begin{equation}
  \tag{12.3.3.4}
  \label{IV.12.3.3.4-fr}
  \operatorname{Ker}(v\otimes1)/\operatorname{Im}(u\otimes1)
  =(\operatorname{Ker}(v)/\operatorname{Im}(u))\otimes_B(B/tB)
\end{equation}
à un isomorphisme canonique près.
\end{lemma}

En effet, l'hypothèse que la suite (12.3.3.2) est exacte entraînera alors
$(\operatorname{Ker}v/\operatorname{Im}u)\otimes_B(B/tB)=0$, et comme $t$
appartient à l'idéal maximal de $B$ et que $\operatorname{Ker}(v)$ est un
$B$-module de type fini (puisque $B$ est supposé noethérien dans (12.3.3.1)),
le lemme de Nakayama prouvera que $\operatorname{Ker}(v)=\operatorname{Im}(u)$.

Reste donc à démontrer (12.3.3.3). Posons $u'=u\otimes1$, $v'=v\otimes1$,
$Z=\operatorname{Ker}(v)$, $H=Z/\operatorname{Im}(u)$, de sorte qu'on a les
suites exactes
\begin{gather*}
  0\to M\to Z\to H\to0,\\
  0\to Z\to N\xrightarrow{v}P\to0,
\end{gather*}
d'où, en tensorisant par $B/tB$ et utilisant le lemme (3.4.1.4) et le fait que
$t$ est $P$-régulier, les suites exactes
\begin{gather*}
  M/tM\xrightarrow{w'}Z/tZ\to H/tH\to0,\\
  0\to Z/tZ\to N/tN\xrightarrow{v'}P/tP\to0,
\end{gather*}
d'où $\operatorname{Ker}(v')=Z/tZ$, et comme
$\operatorname{Im}(w')=\operatorname{Im}(u')$, on obtient (12.3.3.4).

\medskip
II) Supposons maintenant en outre que $\mathcal L_1$ soit $f$-plat~;
remplaçant en outre éventuellement $X$ par $U$, on peut supposer que $X=U$.
Il s'agit de voir que pour tout $x\in X$, on a
$(\mathcal H_1(\mathcal L_\bullet))_x=0$. Soit $s=f(x)$, et soit
$\mathfrak n$ l'idéal $\mathfrak m_s\mathcal O_{X,x}$ qui est contenu dans
l'idéal maximal de $\mathcal O_{X,x}$~; comme
$(\mathcal H_1(\mathcal L_\bullet))_x$ est un $\mathcal O_{X,x}$-module de type
fini ($X$ étant localement noethérien), il est séparé pour la topologie
$\mathfrak n$-adique ($\mathbf 0_{\mathrm I}$, 7.3.5), donc il suffit de
montrer que son \emph{séparé complété} pour cette topologie est 0. Or, en vertu
de (\textbf{III}, 7.4.7.2), ce séparé complété est égal à
$\varprojlim_k\mathcal H_1(\mathcal L_\bullet\otimes_{\mathcal O_S}
(\mathcal O_{S,s}/\mathfrak m_s^{k+1}))$, et il suffira donc de prouver que
chacun des termes de ce système projectif est nul. Mais cela est vrai par
hypothèse pour $k=0$~; raisonnons donc par récurrence sur $k$. La conclusion
résultera du lemme plus général suivant (qu'on appliquera pour
$A=\mathcal O_{S,s}/\mathfrak m_s^{k+1}$ et $\mathfrak J$ égal à l'idéal
maximal $\mathfrak m_s/\mathfrak m_s^{k+1}$ de $A$)~:

\begin{lemma}[12.3.3.5]
\label{IV.12.3.3.5-fr}
Soient $A$ un anneau, $\mathfrak J$ un idéal nilpotent de $A$,
$L_\bullet:P\xrightarrow{u}Q\xrightarrow{v}R$ un complexe de $A$-modules. On
pose $A_0=A/\mathfrak J$,
$L_\bullet^{(0)}=L_\bullet\otimes_AA_0:
P_0\xrightarrow{u_0}Q_0\xrightarrow{v_0}R_0$, et l'on suppose que
$\operatorname{gr}_{\mathfrak J}^{\bullet}(A)$ est un $A_0$-module plat, et
que $Q$ et $R$ sont des $A$-modules plats. Alors la relation
$\mathrm H_1(L_\bullet^{(0)})=0$ entraîne $\mathrm H_1(L_\bullet)=0$.
\end{lemma}

Posons
$L_\bullet^{(k)}=L_\bullet\otimes_A(A/\mathfrak J^{k+1}):
P_k\xrightarrow{u_k}Q_k\xrightarrow{v_k}R_k$~; comme il existe un entier
$k\geq1$ tel que $L_\bullet^{(k)}=L_\bullet$, il suffira de prouver, par
récurrence sur $k$, que la suite
$P_k\xrightarrow{u_k}Q_k\xrightarrow{v_k}R_k$ est
